\documentclass[12pt]{article}
\usepackage[T1]{fontenc}
\usepackage[utf8]{inputenc}
\usepackage{lmodern}
\usepackage{textcomp}
\usepackage{enumitem}
\usepackage{geometry}
\usepackage{graphicx}
\usepackage{amsmath, amssymb}
\geometry{margin=1in}
\begin{document}
\begin{center}
Business Administration - Undergraduate Level Assessment 11 \
Total Marks: 40 \
Answer all questions.
\end{center}

\section*{Multiple Choice (10 marks)}
\begin{enumerate}[label=\arabic*.]
\item What one of the following is not a key management skill in planning?
\begin{enumerate}[label=(\alph*)]
    \item Conceptual skills
    \item Analytical skills
    \item IT and computing skills
    \item Communication skills
\end{enumerate}
\item Which statement is NOT true of a crisis management team?
\begin{enumerate}[label=(\alph*)]
    \item All members should be trained in media relations.
    \item A member of senior management should be included in the team.
    \item A lawyer should be included in the team.
    \item All members should be trained in group decision making.
\end{enumerate}
\item The purposes for which the information is used is in the public's interest.
\begin{enumerate}[label=(\alph*)]
    \item 1,2
    \item 1,3
    \item 2,3
    \item 1,2,3
\end{enumerate}
\item Secondary data has two important advantages over primary data. It is:
\begin{enumerate}[label=(\alph*)]
    \item Capable of compensating for rapid environmental changes and technical improvements.
    \item Always available and complete.
    \item Seldom obsolete and usually fits the dimensions of your problem.
    \item Generally cheaper to gather than primary data and takes less time to find.
\end{enumerate}
\item The process of identifying publics who are involved and affected by a situation central to an organization is called a(n)
\begin{enumerate}[label=(\alph*)]
    \item exploratory survey
    \item situation interview
    \item communication audit
    \item stakeholder analysis
\end{enumerate}
\end{enumerate}

\section*{True / False (10 marks)}
\textit{Answer True or False}
\begin{enumerate}[label=\arabic*.]
\setcounter{enumi}{5}
\item True or False: The Audit Bureau of Circulation does not provide information on readership of publications.
\item True or False: A corporation is independent from its managers, employees, investors, and customers, and it has perpetual succession and owns its own assets.
\item True or False: Subjection to emotional exploitation is experienced through scenarios 1, 3, and 4.
\item True or False: A public issues campaign involves emotionally charged debates between opposing sides that affect their lives.
\item True or False: Services are characterized by intangibility, value, variability, inseparability, and profitability.
\end{enumerate}

\section*{Long Form Response (20 marks)}
\textit{Answer each question with a clear and complete explanation.}
\begin{enumerate}[label=\arabic*.]
\setcounter{enumi}{10}
\item Discuss how the principle of "the greatest good for the greatest number" in utilitarianism applies to decision-making in business and its implications for ethical practices.
\item Discuss the various means by which individuals or organizations can influence government decision-making, focusing on the avenues of approach, breadth of transmission, and content of communication.
\end{enumerate}

\vfill
\noindent\textit{End of Paper}
\end{document}